The release of ChatGPT on 30 November 2022 marked a phase transition in artificial intelligence. A technique called Reinforcement Learning from Human Feedback (RLHF) drove the transition. Christiano et al.~introduced RLHF for Atari trajectory comparisons in 2017. Stiennon et al.~scaled it to text summarization in 2020. Ouyang et al.~then applied it to GPT-3 175B in 2022. The result turned a 175B autoregressive transformer into a usable assistant for hundreds of millions of users. The technical core of that transformation is what this survey addresses: the \emph{alignment} of large language models (LLMs). Gabriel (2020) and Ji et al.~(2023) define alignment as the problem of ensuring that an LLM's outputs are consistent with the intentions, preferences, and values of the humans it serves. Alignment in the LLM era is not an abstract philosophical concern. It is a daily engineering practice that determines whether a model will refuse a malicious request, admit ignorance, follow a JSON schema, or hallucinate a citation. This survey synthesizes the rapidly evolving methods, datasets, benchmarks, and failure modes that have crystallized around this practice between 2017 and 2026, with particular emphasis on Reinforcement Learning from Human Feedback, Direct Preference Optimization, Constitutional AI, scalable oversight, deliberative alignment of reasoning models, and the open problems that remain on the road to systems that may eventually surpass human capabilities. The need for alig\ldots{}
